%%%%%
%
\section{Passage d'un vecteur du rep\`ere g\'eographique dans un rep\`ere mobile}

\itemstyle{Op\'erateur~:} \operfmt{Geo2mobile\_3D}
\index{Geo2mobile\_3D}

%%%%%OBJET%%%%
\itemstyle{Objet~:}

Cet op\'erateur transforme un vecteur 3D donn\'e dans le rep\`ere g�ographique 
dans un rep\`ere mobile caract\'eris\'e par 3 angles de rotation.

%%%%%ARGUMENTS%%%%
\bigskip
%%%%%ENTREES%%%%
\begin{tabarg}
\multicolumn{3}{l}{\titretab {Entr\'ees}}\\
\hline
$X_e$ & Vecteur, \texttt{float} & Composante zonale du vecteur d'entr\'ee, positive vers l'est.\\
$X_n$ & Vecteur, \texttt{float} & Composante m\'eridienne du vecteur d'entr\'ee, positive vers le nord.\\
$X_v$ & Vecteur, \texttt{float} & Composante verticale du vecteur d'entr\'ee, positive vers le haut.\\
$\psi$ & Vecteur, \texttt{float} & Cap de la plateforme mobile ($degres$, dans le sens indirect depuis le nord g�ographique.)\\
$\phi$ & Vecteur, \texttt{float} & Roulis de la plateforme mobile ($degres$, positif lorsque le cot� tribord est enfonc�.)\\
$\theta$ & Vecteur, \texttt{float} & Tangage de la plateforme mobile ($degres$, positif lorsque l'avant est lev�.)\\
\end{tabarg}

\medskip
%%%%%SORTIES%%%%
\begin{tabarg}
\multicolumn{3}{l}{\titretab {Sorties}}\\
\hline
$X_l$ & Vecteur, \texttt{float} & Composante du vecteur r\'esultat suivant l'axe longitudinal de
la plateforme, positif vers l'avant\\
$X_t$ & Vecteur, \texttt{float} & Composante du vecteur r\'esultat suivant l'axe transversal
de la plateforme, positif vers tribord\\
$X_z$ & Vecteur, \texttt{float} & Composante du vecteur r\'esultat suivant l'axe vertical
de la plateforme, positif vers le bas\\
\end{tabarg}


%%%%%FORMULE%%%%
\itemstyle{Formule}

%%Geo2mobile_3D

$X_l =  \cos\theta.\sin\psi . X_e +$
       $\cos\theta.\cos\psi. X_n +$
       $\sin\theta. X_v$ \\

$X_t = (\sin\phi. \sin\theta.\sin\psi + \cos \phi.\cos\psi) . X_e +$
      $(\sin\phi. \sin\theta.\cos\psi - \cos \phi.\sin\psi) . X_n -$
      $\sin\phi.\cos\theta . X_v$ \\

$X_z = (\cos\phi. \sin\theta.\sin\psi - \sin\phi.\cos\psi) . X_e +$
      $(\cos\phi. \sin\theta.\cos\psi + \sin\phi.\sin\psi) . X_n -$
      $\cos\phi.\cos\theta . X_v$ \\



%%%%%AUTEUR%%%%
\itemstyle{Auteur~:}
CNRM/GMEI/TRAMM

%%%%%REFERENCE%%%%
\itemstyle{R\'ef\'erences~:}

D�finition des angles d'attitude dans la norme AFNOR X 02-115 :
\textit{Symboles et vocabulaire de la m�canique du vol}, page 7
\cite{AFNOR-74}.

D�tail des calculs dans le bulletin NCAR-RAF Nr 23 \cite{RAF-23}.
